
来计算下图中
两棵树 $\Gamma_1$、$\Gamma_2$
的情形。
\par\medskip\noindent
\begin{center}
\setlength{\unitlength}{0.75mm}
\begin{picture}(91,47)
\linethickness{0.55pt}
\put(25,39){\line(0,-1){32}}
\put(25,30){\circle*{1.5}}
\put(25,15){\circle*{1.5}}
\put(27,30){\line(2,1){7}}
\put(27,15){\line(2,0){7}}
\put(35,31){$\scriptstyle Y_2$}
\put(35,14){$\scriptstyle Y_1$}
\put(29,21){$\scriptstyle\mathfrak G$}
\put(22,1){$\Gamma_1$}
\put(65,39){\line(0,-1){32}}
\put(65,30){\circle*{1.5}}
\put(65,15){\circle*{1.5}}
\put(67,30){\line(2,1){7}}
\put(67,15){\line(2,0){7}}
\put(75,31){$\scriptstyle Y_1$}
\put(75,14){$\scriptstyle Y_2$}
\put(69,21){$\scriptstyle\mathfrak G$}
\put(62,1){$\Gamma_2$}
\end{picture}
\par\smallskip
{\renewcommand{\thefigure}{7.4.14}\refstepcounter{figure}
\normalfont\bfseries 图~\thefigure\label{fig:7-4-14}}
\end{center}
\par\medskip

注意，
\[
\begin{aligned}
\mathfrak g(\Gamma_1)(x)
&=\Pi_H\bigl(
Y_1(\mathfrak G(Y_2(x)))
\bigr),\\
\mathfrak g(\Gamma_2)(x)
&=(-1)^{\deg'Y_1\deg'Y_2}
\Pi_H\bigl(
Y_2(\mathfrak G(Y_1(x)))
\bigr).
\end{aligned}
\]
计算得
\[
\begin{aligned}
&\overline{\mathfrak m}_1
\bigl(\mathfrak g(\Gamma_1)(x)\bigr)
-(-1)^{\deg Y_1+\deg Y_2+1}
\mathfrak g(\Gamma_1)
\bigl(\overline{\mathfrak m}_1(x)\bigr)\\
&\quad=
\Bigl(
\Pi_H\circ
\bigl(
\overline{\mathfrak m}_1\circ Y_1
-(-1)^{\deg Y_1}
Y_1\circ\overline{\mathfrak m}_1
\bigr)
\circ\mathfrak G\circ Y_2
\Bigr)(x)\\
&\qquad+
(-1)^{\deg'Y_1}
\bigl(
\Pi_H\circ Y_1\circ
(\operatorname{id}-\Pi_H)
\circ Y_2
\bigr)(x)\\
&\qquad+
(-1)^{\deg'Y_1+1}
\Bigl(
\Pi_H\circ Y_1\circ\mathfrak G
\circ
\bigl(
\overline{\mathfrak m}_1\circ Y_2
-(-1)^{\deg Y_2}
Y_2\circ\overline{\mathfrak m}_1
\bigr)
\Bigr)(x).
\end{aligned}
\]
右端第二项包含
\begin{equation}
(-1)^{\deg'Y_1}
\bigl(
\Pi_H\circ Y_1\circ Y_2
\bigr)(x).
\tag{7.4.83.1}\label{eq:7-4-83-1}
\end{equation}
类似地，
\[
\overline{\mathfrak m}_1
\bigl(\mathfrak g(\Gamma_2)(x)\bigr)
-(-1)^{\deg Y_1+\deg Y_2+1}
\mathfrak g(\Gamma_2)
\bigl(\overline{\mathfrak m}_1(x)\bigr)
\]
包含
\begin{equation}
(-1)^{
\deg'Y_2+\deg'Y_1\deg'Y_2
}
\bigl(
\Pi_H\circ Y_2\circ Y_1
\bigr)(x).
\tag{7.4.83.2}\label{eq:7-4-83-2}
\end{equation}
式~\eqref{eq:7-4-83-1}
与~\eqref{eq:7-4-83-2}
之和为
\begin{equation}
\begin{aligned}
&(-1)^{\deg'Y_1}
\Pi_H\circ
\Bigl(
Y_1\circ Y_2
-(-1)^{\deg Y_1\deg Y_2}
Y_2\circ Y_1
\Bigr)(x)\\
&\qquad=
2(-1)^{\deg'Y_1+1}
\bigl(
\Pi_H\circ[Y_1,Y_2]
\bigr)(x).
\end{aligned}
\tag{7.4.83.3}\label{eq:7-4-83-3}
\end{equation}
它与式~\eqref{eq:7-4-82-3}
只差一个整体负号。

下面讨论
式~\eqref{eq:7-4-57}
右端的符号
$*(\mathbf Y_{(1)},\mathbf Y_{(2)},c)$。
如前所述，
式~\eqref{eq:7-4-57}
右端的各项来自
在 $\partial_e Z_\Gamma^{\mathfrak g}$
中出现的
\[
\bigl[
\overline{\mathfrak m}_1,\mathfrak G
\bigr]
=-\operatorname{id}+\Pi_H
\]
的第二项 $\Pi_H$。
按照约定，
须将
$\overline{\mathfrak m}_1$
与属于
$\mathbf Y_{(1)}$
的所有
$\mathfrak G\circ Y_i$
换序，
得到符号
$(-1)^{\deg'\mathbf Y_{(1)}}$。
把这一项移到等式右端，
还产生一个负号；
因此，
$*(\mathbf Y_{(1)},\mathbf Y_{(2)},c)$
中得到 $+1$。
该符号中的其余项，
由前述过程给出。
至此，
命题~\ref{prop:7-4-56}
得证。\hfill$\square$

现在修改
$\mathfrak g_{\ell,k}$
以得到
$\mathfrak g_{\ell,k}^{\circ}$。
定义
\begin{equation}
\mathfrak g_{\ell,k}^{\circ}
(Y_1,\ldots,Y_\ell)(\mathbf x)
=
\mathfrak g_{\ell,k}
(Y_\ell,\ldots,Y_1)(\mathbf x).
\tag{7.4.84}\label{eq:7-4-84}
\end{equation}
